\section{Owner energy at the rank threshold}
\label{sec:actual-harvested-owner-energy}

The energy estimate used in the corrected block alphabet concerns
the actual fine flag threshold. It also subtracts the affine squares
selected by the Gaussian construction. We verify these two changes
to the seed estimate explicitly.

Write $y=r(x)$ for the recording parameter,
$\nu=\mathrm{fiery}(y)$ and
$\tau=\mathrm{phase}(\beta_F(x))$ for the fine rank threshold.
The retiming comparison gives, eventually,
\begin{equation}\label{eq:owner-threshold-band}
 0<\nu,\quad 0<\tau,\qquad \nu/2\le\tau\le2\nu.
\end{equation}
The stores and seed counts remain those of the recorded configuration
at $y$. Only the fine energy rows are read at threshold $\tau$;
the other rows retain their recorded thresholds. In particular,
this operation does not replace the stored flags by flags at $\tau$.

\subsection{Paying the same seeds after changing the fine threshold}

Let $\mathcal F$ be the original fine seed inventory and
$r_i$ the fine plaquette chords. For a threshold $s$, write
$\mathcal F_s=\{i\in\mathcal F:r_i>s\}$. At any
$s\in[\nu/2,2\nu]$, the argument of
Section~\ref{sec:seed-energy} gives
\begin{equation}\label{eq:owner-reread-fine}
 F_s\ge b s\sum_{i\in\mathcal F_s}r_i
            +d s^2|\mathcal F\setminus\mathcal F_s|,
 \qquad b=\lambda_f/2,\quad d=a_f/512.
\end{equation}
Here the guide has norm at most $\epsilon s$, where
$\epsilon=2C/L^2$. The factor two comes from
$\nu\le2s$. To verify the changed constant, a listed small
plaquette has $r_i\ge\nu/4\ge s/8$. If
$\epsilon\le\sqrt{a_f}/16$, then
\[
 \tfrac12(\sqrt{a_f}r_i-\|v_i\|)^2
       \ge a_fs^2/512.
\]
Impose also $\epsilon\sqrt{a_f}\le\lambda_f/4$.
The large-plaquette direct bound then remains
$(3\lambda_f/4)sr_i$. The discarded-path argument of
\eqref{eq:seed-discarded-count} uses witnesses with
$s\exp(\operatorname{dist}/R)<r_j$ and the same original
fine lattice. The recorded clean-neighbourhood conditions,
$s\ge\nu/2$ and the paid path-range bound place each such
witness in $\mathcal F_s$. Thus the discarded rows cost at most
$(\lambda_f/4)s\sum_{\mathcal F_s}r_i$, provided
$\epsilon(1+C_R)C\le\lambda_f/4$. These two estimates prove
\eqref{eq:owner-reread-fine}. No continuity in the clipping
threshold is used.

The next elementary allocation compares the two large-defect
inventories without identifying them. Let $T\ge\nu$ be the
very-large threshold, and let $\mathcal K$ be the fine defects
spent at that threshold. Put
\[
 N=\mathcal F_s,\quad n_0=|\mathcal F\setminus N|,
 \quad R=\sum_{i\in N}r_i,\quad
 V=\sum_{i\in\mathcal F_\nu}r_i.
\]
Then
\[
 V\le R+2\nu n_0,\qquad
 \nu|N|\le2R,\qquad T|\mathcal K|\le V.
\]
The first inequality bounds every term outside $N$ by
$s\le2\nu$; the second uses $s\ge\nu/2$; and the last
uses $\mathcal K\subseteq\mathcal F_\nu$ and
$r_i\ge T$ on $\mathcal K$. Consequently, if $m,c\ge0$
and
\[
 2(3m+c)\le b,\qquad m+2(m+c)\le d/4,
\]
then
\begin{equation}\label{eq:owner-reread-reserve}
 m\nu^2|\mathcal F|+m\nu T|\mathcal K|+c\nu V
       \le F_s.
\end{equation}
Indeed the left side is at most
$(3m+c)\nu R+(3m+2c)\nu^2n_0$.
The displayed conditions and $\nu\le2s$ bound these terms by
$bsR$ and $ds^2n_0$, respectively.

All predictor and coarse rows are unchanged. Apply their reserves
and the hit/unhit comparisons from
Sections~\ref{sec:seed-energy}, \ref{sec:seed-hit-comparison}
and \ref{sec:seed-unhit-comparison}, using
\eqref{eq:owner-reread-reserve} for the fine reserve.
One can choose
\[
 0<m\le\min\{p/32,a_c/64,\lambda_f/32,d/32\}
\]
before increasing $L$ until $c\le m$ and the guide conditions
above hold. Choose the predictor width afterward, so that its
linear expense is at most $pW/16$. The two final smallness
inequalities \eqref{eq:seed-final-smallness}, now with threshold
$\nu$, hold eventually. The same cancellation as in the seed
estimate proves
\begin{equation}\label{eq:owner-rank-income}
 \mathcal E_\tau(O)\ge b_0\nu^2\,s(O),\qquad b_0=m/2>0.
\end{equation}
Here $s(O)$ is the unchanged number of seed shares of the owner
store $O$, and $\mathcal E_\tau(O)$ is its assigned energy
with only the fine rows replaced by their $\tau$ versions.
The finite-volume, tree-slice and local seed hypotheses are those
of the cited seed estimates. Their guide, geometry and comparison
bounds are used on this same store and field.

\subsection{Subtracting the selected affine squares}

Put $B=\lceil(\log y)^2\rceil$. Let $Q$ be an annular
window around $O$, and $H$ a subwindow selecting fine and
predictor rows of $Q$. Use the actual logarithmic signal in
the local Gaussian coordinates, and denote by
$\mathcal X_H$ the sum of its selected affine squares.
Section~\ref{sec:packet-annulus} proves that these are the
actual rows, with the same forest, tag and coarse field, and
that
\begin{equation}\label{eq:owner-annulus-loss}
 \mathcal W(Q)-\mathcal X_H\ge-e(y)s(O),
 \qquad e(y)\ge0,\quad y\,e(y)\longrightarrow0.
\end{equation}
Here $e$ is a function of the recording parameter, evaluated at
$y=r(x)$ in this application. Since $r(x)\to\infty$,
its eventual estimates apply for all sufficiently large $x$.
In fact its local error has upper exponent $-3/2$, and the
number of annular rows per share is bounded by a fixed polynomial
in $B$. This gives the asserted decay and also supplies the
local chart representation of every selected row.

The same bound applies when the fine rows of $Q$ are read at
$\tau$. To see why, the annulus has fine chord at most
$\nu/2\le\tau$. The only possible obstruction is the
exponential path-kept test. Its nontrivial range is contained
in radius $B/50$ by the purchased path-counting bound; in that
range the clean-annulus condition gives chord at most $\nu/2$.
Outside it the right side of the test is at least the universal
chord bound $2$. Thus every relevant path is kept at both
thresholds $\nu$ and $\tau$. Both fine rows have exactly the
same quadratic formula there. The other rows never changed.
This proves equality of the two annular energies, including at
threshold equality.

Since $\nu^2\ge\beta_F(y)^{-1}$ and
$\beta_F(y)/y\to1$, \eqref{eq:owner-annulus-loss} implies
$e(y)\le(b_0/2)\nu^2$ eventually. Adding
\eqref{eq:owner-rank-income} therefore gives
\begin{equation}\label{eq:owner-selected-gain}
 \mathcal E_\tau(O)+\mathcal W_\tau(Q)-\mathcal X_H
 \ge\frac{b_0}{2}\nu^2s(O)
 \ge\frac{b_0}{8}\tau^2s(O).
\end{equation}
This estimate holds for every admitted selected window and
actual chart signal, uniformly in the finite period. The
threshold is chosen after its fixed window and chart constants.

\subsection{The harvested ceiling of an actual packet}

We apply \eqref{eq:owner-selected-gain} to the ceiling used in
the block exponent. Fix the full store $A$, an owner packet
$O$, and separation radius $R\ge205B$. Assume the packet
ownership conditions, the actual seed conditions on $(U,V)$,
positive frame likelihood, and nonzero chart doors. Assume also
$\operatorname{span}\le B/50$, the coarse side at least
$100(1+B)$, and the purchased local chart comparison with
width $w$. The chart radius is $w\nu$, is positive and is
smaller than the first logarithmic radius. These hypotheses
hold before the fine field or its coordinates are varied.

For each signed orientation $\sigma$, turn the owner and full
store, the fine and coarse fields, the forest and the tag together.
Let $Q_\sigma$ be the annulus between radii $2B$ and $200B$
of the turned owner. Let $H_\sigma$ consist of the actual
Gaussian rows selected from this annulus. The span bound keeps
each selected row outside the $2B$ core. The packet separation
$R\ge205B$ keeps the annulus away from other owners. Thus
$H_\sigma$ selects a subwindow of $Q_\sigma$ satisfying
exactly the annular hypotheses used above.

The signal in those rows is the actual signed coordinate array,
not an independently chosen logarithm. On tree bonds both its
increment and the fine field are the identity. On frozen bonds,
the clean chart has norm at most $K\nu$ with $1\le K\le w$;
the cutoff logarithm agrees with that chart because it lies
within the chosen radius. On active bonds the nonzero doors
bound the actual coordinate by $2w\nu$. Therefore this same
array reconstructs the turned fine field and has norm at most
$2w\nu$ throughout every selected stencil. The chart argument
in Section~\ref{sec:packet-annulus} applies with width $2w$.
It proves the hypotheses of \eqref{eq:owner-selected-gain}
for every $H_\sigma$.

Let $\mathcal C_{200B}(O)$ be the oriented elementary energy
addresses whose coarse anchors lie within distance $200B$ of
$O$. The definition of the actual ceiling charge gives the exact
identity
\begin{equation}\label{eq:owner-ceiling-identity}
\begin{split}
 \mathrm{Aff.weight}\sum_{j\in\mathcal C_{200B}(O)}
             \mathrm{levy}_j
 =\operatorname{Av}_{\sigma}\bigl[
       \mathcal W_\tau(\operatorname{anchor}_{200B}(\sigma O))
       -\mathcal X_{H_\sigma}(a_\sigma)
       +\mathcal C_\sigma(m)\bigr].
\end{split}
\end{equation}
Here $a_\sigma$ is the actual array just described, $m$ is
the localized Gaussian mean, and $\mathcal C_\sigma(m)$
is the sum of its charged fine and predictor squares. It is
nonnegative. The average is the normalized signed average;
the factor $\mathrm{Aff.weight}$ has already performed this
normalization.

For completeness, the identity is a finite reindexing. Partition
the oriented energy addresses into fine faces, predictor bonds,
coarse faces and normalization rows. Their anchored sums are
the corresponding four terms of $\mathcal W_\tau$. At a
selected fine or predictor address the charged coordinate square
equals the extracted affine square of the signed array. This
uses the actual path bound and authorized predictor branch.
At an unselected address both subtracted terms are zero. The
mean squares remain with the positive sign in the definition
of the levy. Summing these identities yields
\eqref{eq:owner-ceiling-identity}.

Finally the anchored energy partitions exactly into its core
and annulus:
\[
 \mathcal W_\tau(\operatorname{anchor}_{200B}(\sigma O))
 =\mathcal E_\tau(\sigma O)+\mathcal W_\tau(Q_\sigma).
\]
This is the finite disjoint partition of each of the four row
inventories; the changed fine threshold is used on both sides.
Apply \eqref{eq:owner-selected-gain}, drop the nonnegative
mean squares, and use $s(\sigma O)=s(O)$. Averaging proves
\begin{equation}\label{eq:owner-actual-heft}
 \frac{b_0}{8}\tau^2s(O)
 \le\mathrm{Aff.weight}\sum_{j\in\mathcal C_{200B}(O)}
               \mathrm{levy}_j.
\end{equation}
This is the harvested owner-energy estimate required in
Section~\ref{sec:actual-corrected-body-compiler}. It holds for
all actual compatible packets, forests, tags and coordinates
with the stated seed, likelihood and door conditions. The
uniform chart, branch and seed estimates used to obtain those
conditions remain the previously stated construction inputs;
no assertion about the separate all-compatible warning factor
is made by this energy calculation.
